%% ============================================================
%%  Supplementary Information — standalone document
%%  Target journal: Nature Communications
%%  Build: pdflatex supplementary && bibtex supplementary
%%         && pdflatex supplementary && pdflatex supplementary
%%  Contents: derivation, data provenance, and reproduction instructions
%%  (sourced from section/900_appendix.tex).
%% ============================================================

\documentclass[Numbered]{sn-jnl}

\usepackage[numbers,sort&compress]{natbib}
\usepackage{amsmath,amssymb,amsfonts}
\usepackage{graphicx}
\usepackage{xcolor}
\usepackage{hyperref}
\hypersetup{colorlinks=true, linkcolor=blue, citecolor=blue, urlcolor=blue}
\usepackage{booktabs}
\usepackage{tabularx}
\usepackage{makecell}
\usepackage{pifont}
\usepackage{float}
\usepackage{placeins}
\usepackage{diagbox}

%% ---- Supplementary numbering ------------------------------
%% Sections appear as "Supplementary Note N"; tables as
%% "Supplementary Table N"; equations as (SN) to avoid collision
%% with main-text equation numbers.
\renewcommand{\thesection}{Supplementary Note~\arabic{section}}
\renewcommand{\thesubsection}{\arabic{section}.\arabic{subsection}}
\renewcommand{\tablename}{Supplementary Table}
\renewcommand{\theequation}{S\arabic{equation}}

\begin{document}

\title{Supplementary Information for ``Two dimensionless ratios define
operational limits in hierarchical-memory inference''}

\author*[1]{\fnm{Geunsik} \sur{Lim}}
\email{leemgs@g.skku.edu}
\affil[1]{\orgdiv{College of Information and Communication Engineering},
  \orgname{Sungkyunkwan University (SKKU)},
  \orgaddress{\city{Suwon}, \country{South Korea}}}

\maketitle

\noindent This document provides derivations, data provenance, reproduction
commands, and the protocol boundary accompanying the main text.
Equation numbers of the form (3) refer to the main text;
equations numbered (S1), (S2), \ldots{} are internal to this document.

\section{Derivation and implementation details}

\subsection{Overlap boundary}
With compulsory transfer volume $D$ and sustained link bandwidth $B$, the
transfer time is $T_{\mathrm{transfer}}=D/B$. Perfect overlap can hide at most
$T_{\mathrm{comp}}$, leaving
\begin{equation}
T_{\mathrm{exposed}}\geq
\max(0,T_{\mathrm{transfer}}-T_{\mathrm{comp}}).
\end{equation}
Thus $R_B=T_{\mathrm{comp}}/T_{\mathrm{transfer}}=1$ is the equality boundary.
This is a definition-derived boundary, not a fitted constant.

The manuscript uses a maximum, rather than a sum, for the service-time lower
bound. This permits perfect overlap among compute, memory service, and transfer
and avoids counting the same data movement twice. Any synchronisation,
contention, or incomplete overlap increases latency above that conservative
bound.

\subsection{Residency convention}
The non-resident fraction under the aggregate-capacity approximation is
$\max(0,1-R_C)$. We label $R_C<0.5$ capacity-limited because more than half of
the working set is then non-resident. The choice is transparent and replaceable;
it does not imply a singularity at 0.5.

\subsection{Complete CPU data}
The machine-readable files
\path{code/results/cpu_probe/cpu_hierarchy_records.jsonl} and
\path{code/results/cpu_probe/cpu_hierarchy_summary.json} contain all 14
pointer-chain points and seven-trial summaries. The files
\path{code/results/colab_probe/records.jsonl} and
\path{code/results/colab_probe/summary.json} contain the layered-matrix CPU measurements, and
\path{code/results/colab_probe/tesla_t4_summary.json} contains the per-point
NVIDIA \GpuDevice{} GPU measurements. Values in
the main text are direct rounded summaries of those files.

\subsection{Reproduction commands}
From the repository root, the CPU measurements can be repeated with
\begin{verbatim}
python code/experiments/cpu_hierarchy_probe.py
python code/experiments/colab_regime_measurement.py \
  --backend numpy --quick
python code/experiments/reproduce_regime_map.py
\end{verbatim}
The GPU measurements are reproduced by running the same harness on a
CUDA device (e.g.\ a free cloud GPU runtime); it auto-detects the accelerator:
\begin{verbatim}
python code/experiments/colab_regime_measurement.py
\end{verbatim}
The \texttt{reproduce\_regime\_map.py} command regenerates an analytical
schematic and does not perform an experiment. Measured latencies vary with
processor, GPU, BLAS/CUDA library, system load, and memory placement; exact
numeric identity is not expected on other machines.

\subsection{Accelerator validation status}
A single-GPU run on an NVIDIA \GpuDevice{} was performed with the harness's
CUDA-event path (Section~\ref{sec:gpu}); the XLA/TPU path exists but was not
exercised. This single-accelerator check is not represented as a multi-platform
campaign. A future confirmatory campaign should retain raw device-event
timestamps, software versions, accelerator identifiers, warm-up policy,
excluded-window reasons, and every prespecified operating point across at least
two accelerator architectures.


\end{document}
